\section{Marginals, the joint law, and the ordering corollary}
\label{sec:distribution}

Proposition~\ref{prop:mixed} separates every pair of bin characters.
Finite Fourier inversion therefore separates bin events, including
smoothness at rational powers of the counting scale.  The factorial-moment
limits from Section~\ref{sec:labels} identify their marginals with the
Dickman law.  We then obtain the fixed-scale joint law through every real
counting endpoint, pass to the moving thresholds in Theorem~\ref{thm:main},
and deduce the comparison corollary from the continuous limiting distribution.

\subsection{Finite Fourier inversion}

Adding back the centering in Proposition~\ref{prop:mixed} gives
\begin{equation}
 \lim_{x\to\infty}\frac1x\sum_{1\le n<x}
       \overline{g_x(n)}f_x(n+1)=\overline{\mu_g}\mu.
 \label{eq:distribution-characters}
\end{equation}
Here and throughout this subsection, $x$ runs through positive integers.
Indeed, the mean of $\overline{g_x(n)}$ tends to
$\overline{\mu_g}$ by Lemma~\ref{lem:labels-marginal} and its
initial-segment extension.
For $m\le x$, write
\[
 \boldsymbol b_x(m)
   =\bigl(\Omega_{\mathcal B_{k,x}}(m)\bigr)_{1\le k<J}.
\]
Every coordinate lies in $\{0,\ldots,J\}$: each counted prime
factor exceeds $x^{1/J}$.  Thus reduction modulo $J+1$ is injective
on the set of count vectors.  Put
$\zeta_0=\exp(2\pi i/(J+1))$.  For
$\boldsymbol r\in\{0,\ldots,J\}^{J-1}$ and such a count vector
$\boldsymbol b$, the exact identity
\[
 \ind_{\{\boldsymbol b=\boldsymbol r\}}
  =\frac1{(J+1)^{J-1}}
    \sum_{\boldsymbol h\in\{0,\ldots,J\}^{J-1}}
       \zeta_0^{\boldsymbol h\cdot(\boldsymbol b-\boldsymbol r)}
\]
expresses each bin event as a finite linear combination of characters.
The two phase vectors in \eqref{eq:distribution-characters} are
independent choices, and complex conjugation simply permutes the
available root-of-unity phases.  Applying this identity at $n$ and
$n+1$ proves factorization of the limiting expectations of every
pair of fixed functions of their count vectors.  Both marginal limits
exist by Lemma~\ref{lem:labels-marginal}; the shift of the averaging
range changes a bounded marginal average by $O(1/x)$.

If $c,d\in(0,1)$ are rational, choose $J$ so that both are grid
points.  For $m\le x$, the condition $\Pplus(m)\le x^c$ is
exactly the vanishing of all counts in bins with lower endpoint at
least $x^c$.  Hence the joint smoothness event with thresholds
$x^c,x^d$ has a limiting density equal to the product of its two
marginal limits.  The next calculation identifies those limits.

\subsection{The Dickman marginal}

We recover the classical marginal law, originating with Dickman and
developed by Ramaswami and de Bruijn
\cite{Dickman1930,Ramaswami1949,deBruijn1951}, from the factorial moments
already computed for the bin counts.

\begin{lemma}[Dickman marginal]
\label{lem:distribution-dickman}
For every fixed $c\in(0,1)$,
\begin{equation}
 \lim_{x\to\infty}\frac1x
   \#\{1\le m\le x:\Pplus(m)\le x^c\}=\rho(1/c).
 \label{eq:distribution-dickman}
\end{equation}
The limit is through all real $x$.
\end{lemma}

\begin{proof}
First suppose that $c=k_0/J$ for integers $J\ge2$ and $1\le k_0<J$.
For $m\le x$, set
\[
 N_{c,x}(m)=\sum_{k=k_0}^{J-1}\Omega_{\mathcal B_{k,x}}(m).
\]
Then $\Pplus(m)\le x^c$ exactly when $N_{c,x}(m)=0$, and
$0\le N_{c,x}(m)\le J$.  For a nonnegative integer $z$, write
$(z)_h=z(z-1)\cdots(z-h+1)$, with $(z)_0=1$.  The finite
inclusion--exclusion identity is
\[
 \ind_{\{N_{c,x}(m)=0\}}
    =\sum_{h=0}^{J}\frac{(-1)^h}{h!}(N_{c,x}(m))_h.
\]
For each $0\le h\le J$, the falling-factorial multinomial identity gives
\[
 (N_{c,x}(m))_h
 =\sum_{\substack{r_{k_0},\ldots,r_{J-1}\ge0\\
                   r_{k_0}+\cdots+r_{J-1}=h}}
     \frac{h!}{\prod_{k=k_0}^{J-1}r_k!}
     \prod_{k=k_0}^{J-1}
       \bigl(\Omega_{\mathcal B_{k,x}}(m)\bigr)_{r_k}.
\]
Apply the initial-segment form of \eqref{eq:labels-factorial-limit}
to every term, with $r_k=0$ for $k<k_0$.  The multinomial coefficient
counts the assignments of the $h$ ordered variables to the bins.  Summing the resulting
simplex integrals therefore joins those bins into $(c,1]$ and gives
\[
 \lim_{x\to\infty}\frac1x\sum_{1\le m\le x}(N_{c,x}(m))_h
 =\int_{\substack{s_1,\ldots,s_h\in(c,1]\\
                   s_1+\cdots+s_h\le1}}
       \prod_{i=1}^h\frac{ds_i}{s_i}.
\]
For $h=0$, the empty integral is one, and the left-hand limit is also one.
The initial-segment limit used here holds through real $x$, as does the factorial-moment
limit in Section~\ref{sec:labels}.

In this integral put $s_i=cv_i$.  The upper bound on each $v_i$
is then implied by $v_1+\cdots+v_h\le1/c$, and changing the lower
boundary from $v_i>1$ to $v_i\ge1$ does not change the integral.
Averaging the inclusion--exclusion identity shows that the rational
smoothness density is $H(1/c)$, where for $u\ge0$ we put
\begin{equation}
 H(u)=1+\sum_{h\ge1}\frac{(-1)^h}{h!}I_h(u),\qquad
 I_h(u)=\int_{\substack{v_1,\ldots,v_h\ge1\\
                         v_1+\cdots+v_h\le u}}
                  \prod_{i=1}^h\frac{dv_i}{v_i}.
 \label{eq:distribution-H}
\end{equation}
The sum is locally finite because $I_h(u)=0$ when $h>u$.
In particular, $1/c\le J$ ensures that the sum at $u=1/c$ agrees
with the finite inclusion--exclusion sum above.  Put $I_0(u)=1$
for $u\ge0$.

For $h\ge1$ and $y>z\ge1$, symmetry and the identity
$1=(v_1+\cdots+v_h)/(v_1+\cdots+v_h)$ give
\begin{equation}
 I_h(y)-I_h(z)
       =h\int_z^y I_{h-1}(t-1)\frac{dt}{t}.
 \label{eq:distribution-delay}
\end{equation}
To verify this, cancel the last denominator after using symmetry
and set $t=v_1+\cdots+v_h$.  The remaining variables satisfy
$v_1+\cdots+v_{h-1}\le t-1$.  The same formula holds for $h=1$.
Thus the $I_h$ are continuous, $H=1$ on $[0,1]$, and
\[
 uH'(u)=-H(u-1)\qquad(u>1).
\]
Successive integration on the intervals $[k,k+1]$ uniquely
determines a continuous solution from its values on $[0,1]$.
Hence $H=\rho$, proving \eqref{eq:distribution-dickman} for rational $c$.

For an arbitrary real $c\in(0,1)$, choose rational
$c_-<c<c_+$ in $(0,1)$.  The smoothness event at $x^{c_-}$ is
contained in the event at $x^c$, which is contained in the event at
$x^{c_+}$.  The rational limits bound the lower and upper limits
of the middle normalized count by $\rho(1/c_-)$ and $\rho(1/c_+)$.
Letting $c_-,c_+$ approach $c$ and using continuity of $\rho$ at
the finite argument $1/c$ proves \eqref{eq:distribution-dickman}.
\end{proof}

The marginal law also supplies the endpoint behavior of the limiting
distribution.  It gives $0\le\rho(u)\le1$ for $u>1$, and the same
bounds hold at $u=1$ by definition.  The delay equation makes $\rho$
nonincreasing on $[1,\infty)$.  Its limit at infinity must be zero:
if that limit were $L>0$, integrating
$\rho'(u)=-\rho(u-1)/u\le-L/u$ would contradict nonnegativity.
Consequently
\begin{equation}
 D(t)=
 \begin{cases}
  0,&t\le0,\\
  \rho(1/t),&0<t<1,\\
  1,&t\ge1
 \end{cases}
 \label{eq:distribution-D}
\end{equation}
is a continuous distribution function, including at zero and one,
and its probability measure has no atoms.

\subsection{Fixed and moving thresholds}

Combining the Fourier inversion with Lemma~\ref{lem:distribution-dickman}
gives, for every pair of rational $c,d\in(0,1)$,
\begin{equation}
 \lim_{\substack{x\to\infty\\x\in\N}}\frac1x
 \#\{1\le n<x:\Pplus(n)\le x^c,\
                     \Pplus(n+1)\le x^d\}
     =\rho(1/c)\rho(1/d).
 \label{eq:distribution-rational}
\end{equation}
We first extend this fixed-scale law to arbitrary real exponents and
real counting endpoints.  Fix $c,d\in(0,1)$ and rational exponents
$c_-<c<c_+$ and $d_-<d<d_+$ in $(0,1)$.  Given real $X$, put
$x=\lfloor X\rfloor+1$.  Then $x/X\to1$, the integer ranges
$n\le X$ and $n<x$ agree, and for all sufficiently large $X$,
\[
 x^{c_-}\le X^c\le x^{c_+},\qquad
 x^{d_-}\le X^d\le x^{d_+}.
\]
Thus the normalized count over $1\le n\le X$ with thresholds $X^c,X^d$
lies between $x/X$ times the two normalized counts in
\eqref{eq:distribution-rational} at the lower and upper rational
exponents.  Taking lower and upper limits through real $X$, then
letting the rational exponents approach $c,d$, proves
\begin{equation}
 \begin{split}
 &\lim_{X\to\infty}\frac1X
 \#\{2\le n\le X:\Pplus(n)\le X^c,\ \Pplus(n+1)\le X^d\}\\
 &\qquad=D(c)D(d)=\rho(1/c)\rho(1/d)
 \qquad(0<c,d<1).
 \end{split}
 \label{eq:distribution-fixed}
\end{equation}
The omission of $n=1$ changes at most one summand.  In
\eqref{eq:distribution-fixed}, the exponents are fixed and the limit
is through all real $X$.

This also gives the fixed-scale upper-tail independence law discussed
by Erd\H{o}s and Pomerance~\cite[p.~311]{ErdosPomerance1978}:
\begin{equation}
 \begin{split}
 &\lim_{X\to\infty}\frac1X
 \#\{2\le n\le X:\Pplus(n)>X^c,\ \Pplus(n+1)>X^d\}\\
 &\qquad=(1-D(c))(1-D(d))
 \qquad(0\le c,d\le1).
 \end{split}
 \label{eq:distribution-upper-tail}
\end{equation}
For interior exponents this follows by inclusion--exclusion from
\eqref{eq:distribution-fixed} and the two marginal laws.  The marginal
for $n+1$ is the one in Lemma~\ref{lem:distribution-dickman}, since
shifting its counting range changes only $O(1)$ terms.  Replacing either
strict comparison in \eqref{eq:distribution-upper-tail} by a weak one
has the same limit.  Indeed, for fixed $c>0$, equality
$\Pplus(m)=X^c$ is impossible unless $X^c$ is a prime $p$; in that
case every such $m$ is divisible by $p$, giving only
$O(X^{1-c}+1)=o(X)$ possibilities for $m\le X+1$.  At exponent zero,
the upper-tail condition holds for every $m\ge2$, and equality can
occur only at $m=1$.  At exponent one, an upper-tail condition on
$m\le X+1$ has only $O(1)$ possible integers.  These observations
give the endpoint cases of \eqref{eq:distribution-upper-tail} as well.
Its right-hand side is continuous on $[0,1]^2$ by
\eqref{eq:distribution-D}.

\begin{proof}[Proof of Theorem~\ref{thm:main}]
Fix $a,b\in(0,1)$ and choose fixed exponents
$c_-<a<c_+$ and $d_-<b<d_+$ in $(0,1)$.
For each fixed $\varepsilon\in(0,1)$, all sufficiently large
real $X$ satisfy, uniformly for $\varepsilon X\le n\le X$,
\[
 X^{c_-}\le n^a\le X^{c_+},\qquad
 X^{d_-}\le n^b\le X^{d_+}.
\]
On this interval, the fixed-scale event at the lower exponents is
contained in the moving-threshold event of Theorem~\ref{thm:main},
which is contained in the fixed-scale event at the upper exponents.
The discarded initial interval costs at most
$\varepsilon+O(1/X)$ in normalized counting.  The limit through all
real $X$ in \eqref{eq:distribution-fixed} therefore bounds the lower and
upper limits of the desired count by the corresponding products at the
lower and upper exponents, with that error.  Let
$\varepsilon\downarrow0$, then let the exponents approach $a,b$.
Continuity of $D$ gives the claimed product through all real endpoints,
with ordinary, unweighted counting.
\end{proof}

\subsection{The ordering probability}

\begin{proof}[Proof of Corollary~\ref{cor:comparison}]
For $n\ge2$, put
\[
 U_n=\frac{\log\Pplus(n)}{\log n},\qquad
 V_n=\frac{\log\Pplus(n+1)}{\log n}.
\]
The empirical probability law means uniform sampling from
$2\le n\le\lfloor X\rfloor$; its normalization
$1/(\lfloor X\rfloor-1)$ differs from $1/X$ by a factor tending
to one.  Recall the continuous distribution function $D$ from
\eqref{eq:distribution-D}.

Although $0\le U_n\le1$, the second coordinate can slightly
exceed one.  This overshoot vanishes and does not alter the limiting
law: for every $n\ge N\ge2$,
\[
 0\le V_n\le\frac{\log(n+1)}{\log n}
       \le1+\frac1{N\log N}.
\]
Consequently the empirical laws are tight and every limiting law is
supported on $[0,1]^2$.  More explicitly, clamp $V_n$ to
$\widetilde V_n=\min(V_n,1)$.  For $0<a,b<1$, the joint
distribution functions of $(U_n,\widetilde V_n)$ are exactly those
of $(U_n,V_n)$, hence converge by Theorem~\ref{thm:main} to
$D(a)D(b)$.  These interior rectangles determine the limiting
probability law: their masses approach one as $a,b\uparrow1$,
and finite differences give the masses of all interior grid cells.
Approximating continuous functions on $[0,1]^2$ by such grids proves
weak convergence to the product law with marginal distribution function
$D$.  Removing the clamping does not affect this convergence, by the
displayed bound and the vanishing proportion of $n<N$.

Let $U,V$ be independent with distribution function $D$.
Continuity gives $\Prob(U=V)=0$, and exchangeability of this
limiting pair gives
\[
 \Prob(U<V)=\Prob(V<U)=\frac12.
\]
The boundary of the set $\{(u,v):u<v\}$ is the diagonal, which
has zero product mass.  Weak convergence therefore yields
\[
 \lim_{X\to\infty}\frac1X
    \#\{2\le n\le X:U_n<V_n\}=\frac12.
\]
Since both logarithmic sizes have the same positive denominator,
$U_n<V_n$ is exactly $\Pplus(n)<\Pplus(n+1)$.
Applying the same continuity-set argument to $v<u$ gives the
reverse ordering limit as well.
The symmetry used here belongs to the independent limiting law;
no symmetry of the finite consecutive-integer pairs is assumed.
\end{proof}
